\documentclass[11pt,a4paper]{article}
\usepackage[utf8]{inputenc}
\usepackage[T1]{fontenc}
\usepackage{lmodern}
\usepackage[french]{babel}
\usepackage{amsmath,amssymb,mathtools}
\usepackage{icomma}
\usepackage{xcolor}
\usepackage[most]{tcolorbox}
\usepackage{enumitem}
\usepackage[a4paper,margin=2.2cm,headheight=26pt,headsep=10pt]{geometry}
\usepackage{fancyhdr}
\usepackage[hidelinks]{hyperref}
\definecolor{primary}{RGB}{222,49,99}
\definecolor{primarydark}{RGB}{150,22,55}
\definecolor{excol}{RGB}{0,135,90}
\tcbset{boxbase/.style={enhanced,breakable,boxrule=0.9pt,arc=2.5pt,
  left=3mm,right=3mm,top=2.3mm,bottom=2.3mm,fonttitle=\bfseries,coltitle=white}}
\newtcolorbox[auto counter]{corrige}[1][]{boxbase,colback=excol!4,colframe=excol,
  title={\textsc{Corrigé}~\thetcbcounter\ifx\relax#1\relax\relax\else\textmd{ --- \itshape #1}\fi}}
\newcommand{\R}{\mathbb{R}}
\newcommand{\ds}{\displaystyle}
\pagestyle{fancy}\fancyhf{}
\fancyhead[L]{\small\textcolor{primarydark}{\textbf{Thème 3} — Factorisation et résolution}}
\fancyhead[R]{\small\textcolor{primarydark}{\textsc{Corrigé — Facile}}}
\fancyfoot[C]{\small\thepage}
\renewcommand{\headrulewidth}{0.8pt}
\begin{document}

\begin{center}
{\LARGE\bfseries\color{primarydark} Niveau Facile — Corrigé}\\[2pt]
{\color{primary}\rule{0.45\textwidth}{1.2pt}}
\end{center}
\vspace{1mm}

\begin{corrige}[mise en évidence]
\[ \text{a) }x(2x-3);\quad \text{b) }3x^{2}(2x+3);\quad \text{c) }5x(x+2);\quad
   \text{d) }4x(x^{2}-2);\quad \text{e) }xy(x+y). \]
\end{corrige}

\begin{corrige}[différence de carrés]
\[ \text{a) }(x-3)(x+3);\ \text{b) }(2x-5)(2x+5);\ \text{c) }(x-1)(x+1);\ \text{d) }(3-x)(3+x);\ \text{e) }(4x-1)(4x+1). \]
\end{corrige}

\begin{corrige}[carré parfait]
\[ \text{a) }(x+3)^{2};\quad \text{b) }(x-2)^{2};\quad \text{c) }(3x+1)^{2};\quad
   \text{d) }(x-5)^{2};\quad \text{e) }(2x+1)^{2}. \]
\end{corrige}

\begin{corrige}[discriminant et racines]
\begin{itemize}[leftmargin=1.4em,topsep=1pt,itemsep=1pt]
  \item a) $\Delta=9-8=1>0$, deux racines : $x=\tfrac{3\pm1}{2}$, soit $x=2$ et $x=1$.
  \item b) $\Delta=25-24=1>0$, deux racines : $x=\tfrac{5\pm1}{2}$, soit $x=3$ et $x=2$.
  \item c) $\Delta=4-4=0$, racine double $x=\tfrac{2}{2}=1$.
  \item d) $\Delta=1-4=-3<0$, aucune racine réelle.
\end{itemize}
\end{corrige}

\begin{corrige}[produit nul]
\begin{align*}
&\text{a) }x=1\ \text{ou}\ x=-3; &&\text{b) }x=0\ \text{ou}\ x=\tfrac12; &&\text{c) }x=2\ \text{(double)};\\
&\text{d) }x=-\tfrac52\ \text{ou}\ x=4; &&\text{e) }x=0\ \text{ou}\ x=-7. &&
\end{align*}
\end{corrige}

\end{document}
